% from ega2-1-fr.tex
\oldpage[II]{195}
isomorphismes (\hyperref[II.8.12.5.4-fr]{8.12.5.4}) et
(\hyperref[II.8.12.6.1-fr]{8.12.6.1}), on est ramené à démontrer la
proposition correspondante pour l'homomorphisme canonique
$p_X^*(\operatorname{Proj}_0(\mathcal{M}_X^{\geq}))
\to(\mathcal{M}_X^{\geq})^{\square}|\widehat{E}_X$, autrement dit, on est
ramené au cas où $\mathcal{S}_1$ est un $\mathcal{O}_Y$-Module inversible et
où $\mathcal{S}$ est engendrée par $\mathcal{S}_1$. Avec les notations de
(\hyperref[II.8.12.3-fr]{8.12.3}), on a alors, pour un $f\in S_1$,
$S_f^{\leq}=S_{(f)}[1/f]$ et l'homomorphisme canonique
$M_{(f)}\otimes_{S_{(f)}}S_f^{\leq}\to M_f^{\leq}$ est un isomorphisme par
définition de $M_f^{\leq}$.
\end{proof}

\begin{env}[8.12.9]
\label{II.8.12.9-fr}
Considérons maintenant les $\mathcal{S}$-Modules quasi-cohérents
\[
  \mathcal{M}_{[n]}=\bigoplus_{m\geqslant n}\mathcal{M}_m
\]
et (avec les notations de (\hyperref[II.8.7.2-fr]{8.7.2})) le
$\mathcal{S}^{\natural}$-Module gradué quasi-cohérent
\begin{equation}
\label{II.8.12.9.1-fr}
  \mathcal{M}^{\natural}
  =\left(\bigoplus_{n\geqslant0}\mathcal{M}_{[n]}\right)^{\sim}.
\tag{8.12.9.1}
\end{equation}
On a vu (\hyperref[II.8.7.3-fr]{8.7.3}) qu'il existe un $C$-isomorphisme
canonique $h:C_X\xrightarrow{\sim}\operatorname{Proj}(\mathcal{S}^{\natural})$.
En outre~:
\end{env}

\begin{proposition}[8.12.10]
\label{II.8.12.10-fr}
Il existe un $h$-isomorphisme canonique
\begin{equation}
\label{II.8.12.10.1-fr}
  \operatorname{Proj}_0(\mathcal{M}^{\natural})
  \xrightarrow{\sim}\widetilde{\mathcal{M}}_X.
\tag{8.12.10.1}
\end{equation}
\end{proposition}

\begin{proof}
On raisonne comme dans (\hyperref[II.8.7.3-fr]{8.7.3}), en utilisant cette
fois l'existence du di-isomorphisme (\hyperref[II.8.2.9.3-fr]{8.2.9.3}) au
lieu de (\hyperref[II.8.2.7.3-fr]{8.2.7.3}). Nous laissons les détails au
lecteur.
\end{proof}

\subsection{Fermetures projectives de sous-faisceaux et de sous-schémas
fermés}
\label{subsection:II.8.13-fr}

\begin{env}[8.13.1]
\label{II.8.13.1-fr}
Les hypothèses et notations étant celles de
(\hyperref[II.8.12.1-fr]{8.12.1}), considérons un sous-$\mathcal{S}$-Module
quasi-cohérent $\mathcal{N}$ de $\mathcal{M}$, \emph{non nécessairement
gradué}. On peut donc considérer le $\mathcal{O}_C$-Module quasi-cohérent
$\widetilde{\mathcal{N}}$ associé à $\mathcal{N}$, qui est un
sous-$\mathcal{O}_C$-Module de $\widetilde{\mathcal{M}}$. On a vu d'autre
part (\hyperref[II.8.12.2.1-fr]{8.12.2.1}) que
$\widetilde{\mathcal{M}}$ s'identifie à la restriction de
$\mathcal{M}^{\square}$ à $C$. Comme l'injection canonique
$i:C\to\widehat{C}$ est un morphisme affine
(\hyperref[II.8.3.2-fr]{8.3.2}), et \emph{a fortiori} quasi-compact, le
\emph{prolongement canonique} $(\widetilde{\mathcal{N}})^{-}$, plus grand
sous-$\mathcal{O}_{\widehat{C}}$-Module contenu dans
$\mathcal{M}^{\square}$ et induisant $\widetilde{\mathcal{N}}$ sur $C$, est
un $\mathcal{O}_{\widehat{C}}$-Module quasi-cohérent
(\hyperref[I.9.4.2-fr]{I, 9.4.2}). Nous allons en donner une description
explicite à l'aide d'un $\widehat{\mathcal{S}}$-Module gradué.
\end{env}

\begin{env}[8.13.2]
\label{II.8.13.2-fr}
Pour cela, considérons, pour tout entier $n\geqslant0$, l'homomorphisme
$\bigoplus_{i\leqslant n}\mathcal{M}_i\to\mathcal{M}$ qui, pour tout ouvert
$U$ de $Y$, fait correspondre à la famille
\[
  (s_i)\in\bigoplus_{i\leqslant n}\Gamma(U,\mathcal{M}_i)
\]
la section $\sum_i s_i\in\Gamma(U,\mathcal{M})$. Désignons par
$\mathcal{N}'_n$ l'image réciproque de $\mathcal{N}$ par cet homomorphisme,
qui est un sous-$\mathcal{O}_Y$-Module quasi-cohérent de
$\bigoplus_{i\leqslant n}\mathcal{M}_i$. Considérons maintenant
l'homomorphisme
$\bigoplus_{i\leqslant n}\mathcal{M}_i\to
\widehat{\mathcal{M}}=\mathcal{M}[z]$ qui fait correspondre à $(s_i)$ la
section
$\sum_{i\leqslant n}s_i z^{n-i}\in\Gamma(U,\widehat{\mathcal{M}}_n)$, et
soit $\mathcal{N}_n$ l'image de $\mathcal{N}'_n$ par cet homomorphisme~; on
vérifie aussitôt que
$\overline{\mathcal{N}}=\bigoplus_{n\geqslant0}\mathcal{N}_n$ est un
sous-$\widehat{\mathcal{S}}$-Module gradué (quasi-cohérent) de
$\widehat{\mathcal{M}}$~; on dit que $\overline{\mathcal{N}}$ est déduit de
$\mathcal{N}$ par \emph{homogénéisation}, à l'aide de la «~variable
homogénéisante~» $z$. On notera

